\section*{Anhang A: Detailliertes Komponentenmodell}

Die Anwendung ist als hierarchische Komponentenstruktur in Vue 3 organisiert. Das nachfolgende Diagramm zeigt alle Komponenten, ihre Gruppierung nach Funktion und ihre Abhängigkeiten zur Simulationslogik.

\begin{figure}[htbp]
    \centering
    \resizebox{0.92\textwidth}{!}{
\begin{tikzpicture}[x=0.54cm, y=0.54cm, transform shape,
  comp/.style={draw=black!70, fill=white, align=center, minimum width=1.25cm, minimum height=0.62cm, font=\tiny, rounded corners=1.5pt},
  modal/.style={draw=red!60, fill=red!5, align=center, minimum width=1.25cm, minimum height=0.62cm, font=\tiny, rounded corners=1.5pt},
  visual/.style={draw=green!60, fill=green!5, align=center, minimum width=1.25cm, minimum height=0.62cm, font=\tiny, rounded corners=1.5pt},
  logic/.style={draw=blue!60, fill=blue!5, align=center, minimum width=1.25cm, minimum height=0.62cm, font=\tiny, rounded corners=1.5pt},
  layer/.style={draw=black!40, fill=black!3, rounded corners=2pt, line width=0.5pt},
  arrow/.style={->, >=Latex, line width=0.7pt, black!60},
  title/.style={font=\bfseries\small}
]

% Root
\node[comp] (app) at (0, 16.5) {App.vue};

% AppShell
\node[comp] (appshell) at (0, 15.2) {AppShell.vue};
\draw[arrow] (app) -- (appshell);

% Layer 1: Modal Components
\node[layer, minimum width=15.1cm, minimum height=2.5cm] at (0, 13.3) {};
\node[title, anchor=west] at (-6.1, 14.7) {Modal / Dialog Layer};

\node[modal] (welcome) at (-5.4, 13.6) {Welcome\\Modal};
\node[modal] (algo) at (-3.5, 13.6) {Algorithm\\Picker};
\node[modal] (help) at (-1.6, 13.6) {Knowledge\\Page};
\node[modal] (ranking) at (0.3, 13.6) {Ranking\\Edit};
\node[modal] (menubar) at (2.2, 13.6) {Burger\\Menu};
\node[modal] (radial) at (4.1, 13.6) {Radial\\Menu};
\node[modal] (tour) at (6.0, 13.6) {Tour\\Guide};

\draw[arrow] (appshell) -- (welcome);
\draw[arrow] (appshell) -- (algo);
\draw[arrow] (appshell) -- (help);
\draw[arrow] (appshell) -- (ranking);

% Layer 2: Main Visualization Components
\node[layer, minimum width=15.1cm, minimum height=3.0cm] at (0, 9.2) {};
\node[title, anchor=west] at (-6.1, 11.4) {Visualization / Interaction Layer};

\node[visual] (multiview) at (-5.2, 10.4) {MultiView};
\node[visual] (gantt) at (-3.3, 10.4) {Gantt};
\node[visual] (minigan) at (-1.4, 10.4) {MiniGantt};
\node[visual] (stack) at (0.5, 10.4) {StackList};
\node[visual] (metrics) at (2.4, 10.4) {Metrics\\Table};
\node[visual] (compare) at (4.3, 10.4) {Compare\\Panel};
\node[visual] (ranking-list) at (6.2, 10.4) {Ranking\\List};

\node[visual] (scenario-min) at (-5.2, 8.9) {Scenario\\Min};
\node[visual] (cardstack) at (-3.3, 8.9) {CardStack};
\node[visual] (scrubber) at (-1.4, 8.9) {Scrubber};
\node[visual] (tooltip) at (0.5, 8.9) {Tooltip};

\draw[arrow] (appshell) -- (multiview);
\draw[arrow] (multiview) -- (gantt);
\draw[arrow] (multiview) -- (stack);
\draw[arrow] (multiview) -- (metrics);
\draw[arrow] (multiview) -- (compare);

% Layer 3: Core Logic
\node[layer, minimum width=15.1cm, minimum height=2.6cm] at (0, 4.6) {};
\node[title, anchor=west] at (-6.1, 6.7) {Simulation / Logic Layer};

\node[logic] (simulation) at (-4.0, 5.0) {simulation.ts\\(Core)};
\node[logic] (types) at (-1.3, 5.0) {types.ts\\(Model)};
\node[logic] (compare-util) at (1.4, 5.0) {compare.ts\\(Utils)};
\node[logic] (timeline) at (4.1, 5.0) {timeline\\Layout.ts};

% Composables (Shared Logic)
\node[layer, minimum width=15.1cm, minimum height=2.6cm] at (0, 1.2) {};
\node[title, anchor=west] at (-6.1, 3.1) {Composables (Shared Logic / State Management)};

\node[comp] (workspace) at (-5.7, 1.8) {useScenario\\Workspace};
\node[comp] (playback) at (-3.7, 1.8) {usePlayback};
\node[comp] (comparison) at (-1.7, 1.8) {useComparison};
\node[comp] (gantt-zoom) at (0.3, 1.8) {useGantt\\Zoom};
\node[comp] (timeline-sync) at (2.3, 1.8) {useTimeline\\Sync};
\node[comp] (theme) at (4.3, 1.8) {useTheme};
\node[comp] (focus-trap) at (6.3, 1.8) {useFocus\\Trap};
\node[comp] (tour) at (8.3, 1.8) {useTour};

% Connections
\draw[arrow] (gantt) -- (simulation);
\draw[arrow] (stack) -- (simulation);
\draw[arrow] (metrics) -- (simulation);
\draw[arrow] (compare) -- (compare-util);
\draw[arrow] (gantt) -- (timeline);
\draw[arrow] (minigan) -- (timeline);

\draw[arrow] (multiview) -- (workspace);
\draw[arrow] (gantt) -- (playback);
\draw[arrow] (scrubber) -- (playback);
\draw[arrow] (compare) -- (comparison);

% Legend
\node[anchor=west, font=\tiny, text width=12cm] at (-6.3, -0.4) {
  \textbf{Farbkodierung:} Rot = Modal \quad Grün = Visualisierung \quad Blau = Logik \quad Weiß = Composables
};

\end{tikzpicture}
}


    \caption{Detailliertes Komponentenmodell der Anwendung. Die vier Schichten zeigen: (1) Modal- und Dialog-Komponenten für Eingabe und Einstellungen, (2) Visualisierungs- und Interaktionskomponenten für die Darstellung, (3) Kern-Simulationslogik und Utilities, (4) Shared Composables für Cross-Cutting Concerns wie State Management und Zustandssynchronisation.}
    \label{fig:detailed-component-model}
\end{figure}

\subsection*{Komponentengruppen}

\subsubsection*{Modal/Dialog Layer}
Diese Schicht enthält alle modalen Dialoge und Eingabe-Komponenten:
\begin{itemize}
    \item \textbf{WelcomeModal:} Initiale Begrüßung und Szenario-Auswahl
    \item \textbf{AlgorithmPickerModal:} Auswahl und Konfiguration von Scheduling-Algorithmen
    \item \textbf{KnowledgePage:} Hilfetexte und Wissensdatenbank
    \item \textbf{RankingEditModal:} Konfiguration der Bewertungsmatrix
    \item \textbf{BurgerMenu / RadialMenu:} Navigation und schnelle Aktionen
\end{itemize}

\subsubsection*{Visualization/Interaction Layer}
Die Hauptkomponenten für die Visualisierung und Benutzerinteraktion:
\begin{itemize}
    \item \textbf{MultiView:} Haupt-Container für alle Visualisierungen
    \item \textbf{Gantt:} Gantt-Zeitachsen-Darstellung
    \item \textbf{MiniGantt:} Kompakte Gantt-Darstellung für Vergleichsübersicht
    \item \textbf{StackList:} Darstellung der Warteschlange
    \item \textbf{MetricsTable:} Kennzahlen-Panel
    \item \textbf{ComparisonPanel:} Vergleichsübersicht mehrerer Läufe
    \item \textbf{RankingList:} Bewertungs- und Ranking-Darstellung
\end{itemize}

\subsubsection*{Simulation/Logic Layer}
Kern-Simulationslogik und domänenspezifische Utilities:
\begin{itemize}
    \item \textbf{simulation.ts:} Deterministische Simulationsengine für alle Scheduling-Algorithmen
    \item \textbf{types.ts:} TypeScript-Domänenmodell (Tasks, Events, Metrics, etc.)
    \item \textbf{compare.ts:} Vergleich- und Bewertungs-Utilities
    \item \textbf{timelineLayout.ts:} Layout-Berechnungen für die Gantt-Darstellung
\end{itemize}

\subsubsection*{Composables (Shared Logic)}
Vue 3 Composables für wiederverwandte Logik und Cross-Cutting Concerns:
\begin{itemize}
    \item \textbf{useScenarioWorkspace:} Verwaltung von Szenarien und aktiven Läufen
    \item \textbf{usePlayback:} Kontrolle der Wiedergabe-Logik (Play, Pause, Step)
    \item \textbf{useComparison:} State Management für Vergleichsansicht
    \item \textbf{useGanttZoom:} Zoom- und Pan-Verwaltung für die Gantt-Achse
    \item \textbf{useTimelineSync:} Synchronisation zwischen Gantt, Stack und Kennzahlen
    \item \textbf{useTheme:} Theme- und Dark-Mode-Verwaltung
    \item \textbf{useFocusTrap / useTour:} Accessibility und User Guidance
\end{itemize}

\subsection*{Datenfluss}

Der Datenfluss folgt dem Master-Slave-Prinzip:
\begin{enumerate}
    \item Nutzer-Eingaben in Modalen oder Steuerelementen triggern State-Updates in Composables
    \item Composables orchestrieren die Neuberechnung in der Simulationslogik
    \item simulation.ts produziert neue Ereignisse, Segmente und Metriken
    \item Visualisierungs-Komponenten abonnieren reactive State aus Composables und rendern die neuen Daten
    \item Die Darstellung aktualisiert sich synchron ohne zurückzufließende Berechnungen
\end{enumerate}

Diese unidirektionale Architektur gewährleistet, dass UI und Simulationslogik sich niemals widersprechen.
